%==========================================================
% CHAPTER: Planimetrija
%==========================================================
\chapter{Planimetrija}

\section{Trikampiai}

Trikampis yra viena pagrindinių plokštumos geometrinių figūrų, turinti tris viršūnes,
tris kraštines ir tris kampus. Šiame skyriuje nagrinėjamos esminės trikampio savybės,
ploto formulės ir pagrindinės trigonometrinės teoremos, būtinos VBE II daliai.

% -------------------------------------------------------
\subsection{Trikampio savybės ir požymiai}
% -------------------------------------------------------

\begin{definition}
\textbf{Trikampis} $ABC$ – tai figūra, sudaryta iš trijų taškų $A$, $B$, $C$,
nepriklausančių vienai tiesei, ir trijų atkarpų $AB$, $BC$, $CA$, vadinamų
\emph{kraštinėmis}. Taškai $A$, $B$, $C$ vadinami \emph{viršūnėmis}, o kampai
$\alpha = \angle BAC$, $\beta = \angle ABC$, $\gamma = \angle BCA$ – \emph{vidiniais kampais}.
\end{definition}

\subsubsection*{Pagrindinės trikampio savybės}

\begin{theorem}[Kampų suma]
Bet kurio trikampio vidaus kampų suma lygi $180°$:
\[
\alpha + \beta + \gamma = 180°.
\]
\end{theorem}

\begin{proof}
Tegu $ABC$ – trikampis. Per viršūnę $B$ nubrėžkime tiesę $l$, lygiagrečią kraštinei $AC$.
Lygiagrečioms tiesėms $l$ ir $AC$, kertamoms skersiniais $AB$ ir $BC$,
priešiniai kampai yra lygūs, todėl:
\[
\angle ABl_1 = \alpha, \quad \angle CBl_2 = \gamma.
\]
Tiesia linija sudaromas $180°$ kampas, todėl
$\alpha + \beta + \gamma = 180°$. \qed
\end{proof}

\begin{corollary}
Išorinis trikampio kampas lygus dviejų jam neištisinių vidaus kampų sumai:
\[
\delta = \alpha + \beta.
\]
\end{corollary}

\subsubsection*{Trikampių rūšys}

Pagal \textbf{kampus}:
\begin{itemize}
  \item \textit{Smailakampis} – visi trys kampai smailūs ($< 90°$).
  \item \textit{Stačiakampis} – vienas kampas status ($= 90°$).
  \item \textit{Bukakampis} – vienas kampas bukas ($> 90°$).
\end{itemize}

Pagal \textbf{kraštines}:
\begin{itemize}
  \item \textit{Lygiakraštis} – visos trys kraštinės lygios.
  \item \textit{Lygiašonis} – dvi kraštinės lygios.
  \item \textit{Įvairiakraštis} – visos kraštinės skirtingos.
\end{itemize}

\subsubsection*{Trikampių lygybė}

\begin{definition}
Du trikampiai vadinami \textbf{lygiais}, jei jų atitinkamos kraštinės ir atitinkami kampai
yra lygūs.
\end{definition}

\begin{theorem}[Trikampių lygybės požymiai]
Du trikampiai yra lygūs, jeigu tenkinama viena iš šių sąlygų:
\begin{enumerate}
  \item \textbf{(KKK)} Visos trys atitinkamos kraštinės yra lygios:
        $a = a'$, $b = b'$, $c = c'$.
  \item \textbf{(KKa)} Dvi kraštinės ir tarp jų esantis kampas yra lygūs:
        $a = a'$, $b = b'$, $\gamma = \gamma'$.
  \item \textbf{(KaK)} Viena kraštinė ir prie jos esantys du kampai yra lygūs:
        $c = c'$, $\alpha = \alpha'$, $\beta = \beta'$.
\end{enumerate}
\end{theorem}

\begin{tcolorbox}[colback=blue!5!white, colframe=blue!60!black, title=Stačiakampio trikampio lygybės požymis]
Stačiakampiai trikampiai yra lygūs, jeigu jų statiniai yra lygūs arba
įžambinė ir vienas statinis yra lygūs.
\end{tcolorbox}

\subsubsection*{Trikampių panašumas}

\begin{definition}
Du trikampiai vadinami \textbf{panašiais} ($\triangle ABC \sim \triangle A'B'C'$),
jeigu jų atitinkami kampai yra lygūs ir atitinkamų kraštinių santykiai lygūs:
\[
\frac{a}{a'} = \frac{b}{b'} = \frac{c}{c'} = k,
\]
kur $k$ – \emph{panašumo koeficientas}.
\end{definition}

\begin{theorem}[Trikampių panašumo požymiai]
Du trikampiai yra panašūs, jeigu:
\begin{enumerate}
  \item \textbf{(KK)} Du atitinkami kampai yra lygūs.
  \item \textbf{(KKK)} Trys atitinkamų kraštinių santykiai yra lygūs.
  \item \textbf{(KKa)} Dvi kraštinių poros santykiai lygūs ir tarp jų esantys
        kampai lygūs.
\end{enumerate}
\end{theorem}

\begin{remark}
Panašių figūrų perimetrų santykis lygus panašumo koeficientui $k$,
o plotų santykis lygus $k^2$.
\end{remark}

\subsubsection*{Svarbūs trikampio elementai}

\begin{itemize}
  \item \textbf{Aukštinė} – atkarpa iš viršūnės, statmena priešingai kraštinei.
        Trys aukštinės kertasi viename taške – \emph{ortocentre} $H$.
  \item \textbf{Mediana} – atkarpa, jungianti viršūnę su priešingos kraštinės viduriu.
        Trys medianos kertasi \emph{sunkio centre} $G$, kuris dalija kiekvieną medianą
        santykiu $2:1$ nuo viršūnės.
  \item \textbf{Pusiaukampinė} – spindulys, dalijantis kampą per pusę.
        Trys pusiaukampinės kertasi \emph{įbrėžtinio apskritimo centre} $I$.
  \item \textbf{Vidurio statmuo} – statmuo kraštinės vidurio taške.
        Trys vidurio statmenys kertasi \emph{apibrėžtinio apskritimo centre} $O$.
\end{itemize}

\begin{theorem}[Pitagoro teorema]
Stačiakampiame trikampyje įžambinės ilgio kvadratas lygus abiejų statinių ilgių
kvadratų sumai:
\[
c^2 = a^2 + b^2,
\]
kur $c$ – įžambinė, $a$ ir $b$ – statiniai.
\end{theorem}

\begin{example}
Stačiakampiame trikampyje statiniai yra $a = 3$ cm ir $b = 4$ cm. Rasti įžambinę.
\[
c = \sqrt{a^2 + b^2} = \sqrt{9 + 16} = \sqrt{25} = 5 \text{ cm}.
\]
\end{example}

\begin{exercise}
Trikampio $ABC$ kampai yra $\alpha = 47°$ ir $\beta = 63°$. Raskite trečiąjį kampą $\gamma$.
\end{exercise}

% -------------------------------------------------------
\subsection{Trikampio plotas}
% -------------------------------------------------------

\subsubsection*{Pagrindinės ploto formulės}

Trikampio plotui apskaičiuoti naudojamos kelios formulės, priklausomai nuo žinomų elementų.

\begin{theorem}[Pagrindinė ploto formulė]
Trikampio plotas lygus pusės pagrindo ir atitinkamos aukštinės sandaugai:
\[
S = \frac{1}{2} \cdot a \cdot h_a,
\]
kur $a$ – pasirinkta kraštinė (pagrindas), $h_a$ – aukštinė į tą kraštinę.
\end{theorem}

\begin{theorem}[Plotas per dvi kraštines ir kampą]
Jeigu žinomos dvi trikampio kraštinės $b$, $c$ ir tarp jų esantis kampas $\alpha$, tai:
\[
\boxed{S = \frac{1}{2} \cdot b \cdot c \cdot \sin\alpha.}
\]
\end{theorem}

\begin{proof}
Trikampyje $ABC$ nubrėžkime aukštinę $h$ iš viršūnės $C$ į kraštinę $AB = c$.
Tuomet $h = b \sin\alpha$ (iš stačiakampio trikampio apibrėžimo).
Įstatę į pagrindinę formulę:
\[
S = \frac{1}{2} \cdot c \cdot h = \frac{1}{2} \cdot c \cdot b\sin\alpha = \frac{1}{2}bc\sin\alpha. \quad \qed
\]
\end{proof}

\begin{tcolorbox}[colback=green!5!white, colframe=green!60!black, title=Visos trys kraštinių–kampo formulės]
\[
S = \frac{1}{2} a b \sin\gamma = \frac{1}{2} b c \sin\alpha = \frac{1}{2} a c \sin\beta.
\]
\end{tcolorbox}

\subsubsection*{Herono formulė}

\begin{theorem}[Herono formulė]
Jeigu žinomos visos trys trikampio kraštinės $a$, $b$, $c$, tai plotas:
\[
S = \sqrt{p(p-a)(p-b)(p-c)},
\]
kur $p = \dfrac{a+b+c}{2}$ – \emph{pusperimetris}.
\end{theorem}

\begin{example}
Raskite trikampio, kurio kraštinės $a = 5$, $b = 6$, $c = 7$, plotą.

Pusperimetris: $p = \dfrac{5+6+7}{2} = 9$.
\[
S = \sqrt{9 \cdot (9-5)(9-6)(9-7)} = \sqrt{9 \cdot 4 \cdot 3 \cdot 2} = \sqrt{216} = 6\sqrt{6} \approx 14{,}70 \text{ cm}^2.
\]
\end{example}

\subsubsection*{Plotas per įbrėžtinį ir apibrėžtinį apskritimus}

\begin{theorem}
Jeigu $r$ – įbrėžtinio apskritimo spindulys, o $p$ – pusperimetris, tai:
\[
S = p \cdot r.
\]
\end{theorem}

\begin{theorem}
Jeigu $R$ – apibrėžtinio apskritimo spindulys, tai:
\[
S = \frac{abc}{4R}.
\]
\end{theorem}

\begin{exercise}
Lygiašonio trikampio šoninė kraštinė $b = 10$ cm, pagrindo kampas $\beta = 40°$.
Apskaičiuokite trikampio plotą.
\end{exercise}

\begin{exercise}
Trikampio kraštinės yra $a = 8$, $b = 15$, $c = 17$. Patikrinkite, ar trikampis stačiakampis,
ir apskaičiuokite jo plotą abiem būdais: pagrindinę formule ir Herono formule.
\end{exercise}

% -------------------------------------------------------
\subsection{Sinusų ir kosinusų teoremos}
% -------------------------------------------------------

Sinusų ir kosinusų teoremos leidžia spręsti \emph{bet kurį} trikampį –
t.\,y. rasti nežinomas kraštines ir kampus, žinant pakankamai pradinių duomenų.

\subsubsection*{Sinusų teorema}

\begin{theorem}[Sinusų teorema]
Trikampio kraštinių ilgiai proporcingi prieš jas esančių kampų sinusams, o šis santykis
lygus apibrėžtinio apskritimo skersmeniui $2R$:
\[
\boxed{\frac{a}{\sin\alpha} = \frac{b}{\sin\beta} = \frac{c}{\sin\gamma} = 2R,}
\]
kur $R$ – apibrėžtinio apskritimo spindulys, $a, b, c$ – kraštinės,
$\alpha, \beta, \gamma$ – jiems priešingi kampai.
\end{theorem}

\begin{proof}
Tegu $ABC$ – smailakampis trikampis ir $O$ – jo apibrėžtinio apskritimo centras,
$R$ – spindulys. Nubrėžkime skersmenį $BD$ per viršūnę $B$.

Kampas $\angle BAC = \alpha$ yra įbrėžtinis kampas, o $\angle BDC$ – taip pat
įbrėžtinis kampas, remiasi tą pačią lanką $BC$, todėl $\angle BDC = \alpha$.

Trikampyje $BDC$: $\angle BCD = 90°$ (kaip kampas, įbrėžtas į pusapskritimį), todėl:
\[
\sin\alpha = \sin(\angle BDC) = \frac{BC}{BD} = \frac{a}{2R}
\implies \frac{a}{\sin\alpha} = 2R.
\]
Analogiškai $\dfrac{b}{\sin\beta} = 2R$ ir $\dfrac{c}{\sin\gamma} = 2R$. \qed
\end{proof}

\subsubsection*{Sinusų teoremos taikymas}

Sinusų teorema taikoma, kai žinoma:
\begin{enumerate}
  \item Dvi kraštinės ir vienam iš jų priešingas kampas (\textbf{KKa} atvejis).
  \item Viena kraštinė ir du kampai (\textbf{KaK} arba \textbf{KKa} atvejis).
\end{enumerate}

\begin{example}
Trikampyje $ABC$: $a = 7$ cm, $\alpha = 50°$, $\beta = 70°$.
Raskite kraštinę $b$ ir apibrėžtinio apskritimo spindulį $R$.

\textit{Sprendimas.}
\[
\frac{a}{\sin\alpha} = \frac{b}{\sin\beta}
\implies b = \frac{a \cdot \sin\beta}{\sin\alpha}
= \frac{7 \cdot \sin 70°}{\sin 50°}
= \frac{7 \cdot 0{,}9397}{0{,}7660}
\approx 8{,}58 \text{ cm}.
\]
\[
2R = \frac{a}{\sin\alpha} = \frac{7}{\sin 50°} \approx \frac{7}{0{,}7660} \approx 9{,}14 \text{ cm},
\quad R \approx 4{,}57 \text{ cm}.
\]
\end{example}

\subsubsection*{Kosinusų teorema}

\begin{theorem}[Kosinusų teorema]
Trikampio kraštinės ilgio kvadratas lygus kitų dviejų kraštinių ilgių kvadratų sumai
minus dviguba tų kraštinių ilgių ir tarp jų esančio kampo kosinuso sandauga:
\[
\boxed{a^2 = b^2 + c^2 - 2bc\cos\alpha.}
\]
Analogiškai:
\[
b^2 = a^2 + c^2 - 2ac\cos\beta, \qquad
c^2 = a^2 + b^2 - 2ab\cos\gamma.
\]
\end{theorem}

\begin{proof}
Tegu trikampyje $ABC$ aukštinė $h$ nubrėžta iš viršūnės $C$ į kraštinę $AB$.
Pažymėkime $BD = x$, tada $AD = c - x$. Pagal Pitagoro teoremą:
\[
h^2 = a^2 - x^2 \quad \text{ir} \quad h^2 = b^2 - (c-x)^2.
\]
Iš čia:
\[
a^2 - x^2 = b^2 - c^2 + 2cx - x^2
\implies a^2 = b^2 - c^2 + 2cx + x^2... 
\]
Iš stačiakampio trikampio $ABС$: $x = c - b\cos\alpha$ ... tačiau paprasčiau:
iš $\triangle ABС$ aukštinė $h = b\sin\alpha$, o $BD = b\cos\alpha$, todėl
$AD = c - b\cos\alpha$. Pagal Pitagoro teoremą $\triangle ACD$:
\[
a^2 = h^2 + (c - b\cos\alpha)^2
= b^2\sin^2\alpha + c^2 - 2bc\cos\alpha + b^2\cos^2\alpha.
\]
Kadangi $\sin^2\alpha + \cos^2\alpha = 1$:
\[
a^2 = b^2 + c^2 - 2bc\cos\alpha. \qquad \qed
\]
\end{proof}

\begin{remark}
Kai $\alpha = 90°$, gauname $\cos 90° = 0$, todėl $a^2 = b^2 + c^2$ –
tai \textbf{Pitagoro teorema}. Taigi kosinusų teorema yra Pitagoro teoremos apibendrinimas.
\end{remark}

\subsubsection*{Kosinusų teoremos taikymas}

Kosinusų teorema taikoma, kai žinoma:
\begin{enumerate}
  \item Trys kraštinės (\textbf{KKK} atvejis) – ieškome kampų.
  \item Dvi kraštinės ir tarp jų esantis kampas (\textbf{KKa} atvejis) – ieškome trečiosios kraštinės.
\end{enumerate}

Kampams rasti iš kosinusų teoremos išreiškiame:
\[
\cos\alpha = \frac{b^2 + c^2 - a^2}{2bc}, \quad
\cos\beta = \frac{a^2 + c^2 - b^2}{2ac}, \quad
\cos\gamma = \frac{a^2 + b^2 - c^2}{2ab}.
\]

\begin{example}
Trikampio kraštinės $a = 5$, $b = 7$, $c = 8$. Raskite kampą $\alpha$.
\[
\cos\alpha = \frac{b^2 + c^2 - a^2}{2bc}
= \frac{49 + 64 - 25}{2 \cdot 7 \cdot 8}
= \frac{88}{112}
= \frac{11}{14} \approx 0{,}7857.
\]
\[
\alpha = \arccos(0{,}7857) \approx 38{,}2°.
\]
\end{example}

\begin{example}
Trikampyje $ABC$: $b = 6$ cm, $c = 9$ cm, $\alpha = 120°$. Raskite kraštinę $a$.
\[
a^2 = b^2 + c^2 - 2bc\cos\alpha
= 36 + 81 - 2 \cdot 6 \cdot 9 \cdot \cos 120°
= 117 - 108 \cdot \left(-\frac{1}{2}\right)
= 117 + 54 = 171.
\]
\[
a = \sqrt{171} = 3\sqrt{19} \approx 13{,}08 \text{ cm}.
\]
\end{example}

\subsubsection*{Teoremų taikymo pasirinkimas}

\begin{center}
\begin{tabular}{|c|c|c|}
\hline
\textbf{Žinomi elementai} & \textbf{Taikoma teorema} & \textbf{Atvejis} \\
\hline
Dvi kraštinės + tarp jų kampas & Kosinusų & KKa \\
Trys kraštinės & Kosinusų & KKK \\
Viena kraštinė + du kampai & Sinusų & KaK \\
Dvi kraštinės + priešingas kampas & Sinusų & KKa \\
\hline
\end{tabular}
\end{center}

\begin{exercise}
Orientaciniame žaidime reikia pasiekti tašką $C$. Žinoma, kad $|AB| = 1{,}2$ km,
$\angle BAC = 38°$ ir $\angle ABC = 75°$. Apskaičiuokite atstumus $|AC|$ ir $|BC|$.
\end{exercise}

\begin{exercise}
Lygiagretainio kraštinės $a = 5$ cm ir $b = 8$ cm, o vienas kampas $\varphi = 60°$.
Apskaičiuokite abiejų įstrižainių ilgius.

\textit{Užuomina:} Lygiagretainio įstrižainėse pasinaudokite kosinusų teorema,
atsižvelgdami į tai, kad gretimi lygiagretainio kampai sudaro $180°$.
\end{exercise}

\begin{tcolorbox}[colback=yellow!10!white, colframe=orange!80!black, title=VBE užduočių tipai – trikampiai]
\begin{itemize}
  \item Rasti nežinomą kraštinę arba kampą, taikant sinusų ar kosinusų teoremą.
  \item Apskaičiuoti trikampio plotą, žinant dvi kraštines ir kampą.
  \item Spręsti uždavinius apie apibrėžtinį arba įbrėžtinį apskritimą.
  \item Praktiniai uždaviniai: atstumai, aukščiai, orientavimasis.
\end{itemize}
\end{tcolorbox}

\section{Keturkampiai}

Keturkampis – tai daugiakampis, turintis keturias kraštines, keturias viršūnes ir keturis kampus. Keturkampio kampų suma lygi $360°$. Šiame skyriuje nagrinėjami svarbiausi specialieji keturkampiai: lygiagretainis ir jo atvejai (stačiakampis, rombas, kvadratas) bei trapecija.

% -------------------------------------------------------
\subsection{Lygiagretainis, stačiakampis, rombas, kvadratas}
% -------------------------------------------------------

\begin{definition}[Lygiagretainis]
\textbf{Lygiagretainis} – tai keturkampis, kurio priešingosios kraštinės yra lygiagrečios poromis:
\[
AB \parallel CD \quad \text{ir} \quad BC \parallel AD.
\]
\end{definition}

\subsubsection*{Lygiagretainio savybės}

\begin{theorem}[Lygiagretainio savybės]
Lygiagretainiui $ABCD$ galioja šios savybės:
\begin{enumerate}
    \item Priešingosios kraštinės yra lygios: $AB = CD$, $\; BC = AD$.
    \item Priešingieji kampai yra lygūs: $\angle A = \angle C$, $\; \angle B = \angle D$.
    \item Gretimų kampų suma lygi $180°$: $\angle A + \angle B = 180°$.
    \item Įstrižainės susikerta ir susikirtimo taške vieną kitą dalija pusiau:
    \[
    AO = OC, \quad BO = OD,
    \]
    kur $O$ – įstrižainių susikirtimo taškas.
    \item Įstrižainių kvadratų suma lygi kraštinių kvadratų sumai:
    \[
    d_1^2 + d_2^2 = 2(a^2 + b^2).
    \]
\end{enumerate}
\end{theorem}

\begin{remark}
Norint įrodyti, kad keturkampis yra lygiagretainis, pakanka patikrinti \emph{vieną} iš šių požymių: priešingos kraštinės lygios, priešingi kampai lygūs, įstrižainės dalija viena kitą pusiau, arba dvi priešingos kraštinės yra tuo pačiu metu lygios ir lygiagrečios.
\end{remark}

\subsubsection*{Lygiagretainio plotas ir perimetras}

Lygiagretainio \textbf{aukštinė} $h$ – tai statmuo, nuleistas iš vienos kraštinės į priešingą (arba jos tęsinį).

\begin{tcolorbox}[colback=blue!5!white, colframe=blue!60!black, title=Lygiagretainio ploto formulės]
\[
S = a \cdot h_a = b \cdot h_b = ab\sin\alpha,
\]
kur $a$, $b$ – lygiagretainio kraštinės, $h_a$ – į kraštinę $a$ nuleista aukštinė, $\alpha$ – kampas tarp gretimų kraštinių.
\end{tcolorbox}

\begin{tcolorbox}[colback=green!5!white, colframe=green!60!black, title=Lygiagretainio perimetras]
\[
P = 2(a + b).
\]
\end{tcolorbox}

\begin{example}
Lygiagretainio kraštinės $a = 8$ cm ir $b = 5$ cm, kampas tarp jų $\alpha = 30°$. Raskite lygiagretainio plotą.

\textbf{Sprendimas.}
\[
S = ab\sin\alpha = 8 \cdot 5 \cdot \sin 30° = 40 \cdot \frac{1}{2} = 20 \text{ cm}^2.
\]
\end{example}

\begin{exercise}
Lygiagretainio plotas $S = 36$ cm$^2$, viena kraštinė $a = 9$ cm. Raskite į šią kraštinę nuleistos aukštinės ilgį.
\end{exercise}

% -------------------------------------------------------
\subsubsection*{Stačiakampis}

\begin{definition}[Stačiakampis]
\textbf{Stačiakampis} – tai lygiagretainis, kurio visi kampai yra tiesieji ($90°$).
\end{definition}

Stačiakampis paveldi visas lygiagretainio savybes, tačiau turi ir papildomų:

\begin{theorem}[Stačiakampio savybės]
\begin{enumerate}
    \item Visi keturi kampai lygūs $90°$.
    \item Įstrižainės yra lygios: $d_1 = d_2$.
    \item Įstrižainės susikerta ir dalija viena kitą pusiau, todėl kiekviena iš jų yra lygi pusei kitos.
    \item Įstrižainės ilgis (iš Pitagoro teoremos):
    \[
    d = \sqrt{a^2 + b^2}.
    \]
\end{enumerate}
\end{theorem}

\begin{tcolorbox}[colback=blue!5!white, colframe=blue!60!black, title=Stačiakampio formulės]
\[
S = a \cdot b, \qquad P = 2(a+b), \qquad d = \sqrt{a^2 + b^2}.
\]
\end{tcolorbox}

\begin{example}
Stačiakampio kraštinės $a = 6$ cm ir $b = 8$ cm. Raskite jo įstrižainės ilgį ir plotą.

\textbf{Sprendimas.}
\[
d = \sqrt{6^2 + 8^2} = \sqrt{36 + 64} = \sqrt{100} = 10 \text{ cm}.
\]
\[
S = 6 \cdot 8 = 48 \text{ cm}^2.
\]
\end{example}

\begin{exercise}
Stačiakampio plotas $S = 60$ cm$^2$, o viena kraštinė $a = 12$ cm. Raskite jo perimetrą ir įstrižainės ilgį.
\end{exercise}

% -------------------------------------------------------
\subsubsection*{Rombas}

\begin{definition}[Rombas]
\textbf{Rombas} – tai lygiagretainis, kurio visos keturios kraštinės yra lygios:
\[
AB = BC = CD = DA = a.
\]
\end{definition}

\begin{theorem}[Rombo savybės]
\begin{enumerate}
    \item Visos kraštinės lygios: $a = b$.
    \item Priešingi kampai lygūs: $\angle A = \angle C$, $\; \angle B = \angle D$.
    \item Įstrižainės yra statmenos viena kitai: $d_1 \perp d_2$.
    \item Įstrižainės dalija kampo pusiaukampinėmis pusiau.
    \item Ryšys tarp kraštinės ir įstrižainių (iš Pitagoro teoremos):
    \[
    a^2 = \left(\frac{d_1}{2}\right)^2 + \left(\frac{d_2}{2}\right)^2 \implies d_1^2 + d_2^2 = 4a^2.
    \]
\end{enumerate}
\end{theorem}

\begin{tcolorbox}[colback=blue!5!white, colframe=blue!60!black, title=Rombo ploto formulės]
\[
S = a^2 \sin\alpha = a \cdot h = \frac{d_1 \cdot d_2}{2},
\]
kur $\alpha$ – rombo kampas, $h$ – aukštinė, $d_1$, $d_2$ – įstrižainių ilgiai.
\end{tcolorbox}

\begin{tcolorbox}[colback=green!5!white, colframe=green!60!black, title=Rombo perimetras]
\[
P = 4a.
\]
\end{tcolorbox}

\begin{example}
Rombo įstrižainės $d_1 = 10$ cm ir $d_2 = 24$ cm. Raskite rombo kraštinę ir plotą.

\textbf{Sprendimas.} Kadangi įstrižainės statmenos ir dalija viena kitą pusiau:
\[
a = \sqrt{\left(\frac{d_1}{2}\right)^2 + \left(\frac{d_2}{2}\right)^2} = \sqrt{5^2 + 12^2} = \sqrt{25 + 144} = \sqrt{169} = 13 \text{ cm}.
\]
\[
S = \frac{d_1 \cdot d_2}{2} = \frac{10 \cdot 24}{2} = 120 \text{ cm}^2.
\]
\end{example}

\begin{exercise}
Rombo kraštinė $a = 5$ cm, o vienas jo kampas lygus $120°$. Raskite rombo plotą ir įstrižainių ilgius.
\end{exercise}

% -------------------------------------------------------
\subsubsection*{Kvadratas}

\begin{definition}[Kvadratas]
\textbf{Kvadratas} – tai stačiakampis, kurio visos kraštinės lygios, arba, ekvivalenčiai, rombas, kurio visi kampai tiesieji.
\end{definition}

Kvadratas yra ir stačiakampis, ir rombas vienu metu, todėl pasižymi visomis abiejų figūrų savybėmis:

\begin{theorem}[Kvadrato savybės]
\begin{enumerate}
    \item Visos kraštinės lygios: $AB = BC = CD = DA = a$.
    \item Visi kampai lygūs $90°$.
    \item Įstrižainės lygios, statmenos ir dalija viena kitą pusiau.
    \item Įstrižainės dalija kampus $45°$ kampais.
    \item Įstrižainės ilgis: $d = a\sqrt{2}$.
\end{enumerate}
\end{theorem}

\begin{tcolorbox}[colback=blue!5!white, colframe=blue!60!black, title=Kvadrato formulės]
\[
S = a^2, \qquad P = 4a, \qquad d = a\sqrt{2}.
\]
\end{tcolorbox}

\begin{example}
Kvadrato įstrižainė $d = 6\sqrt{2}$ cm. Raskite jo plotą ir perimetrą.

\textbf{Sprendimas.} Iš $d = a\sqrt{2}$:
\[
a = \frac{d}{\sqrt{2}} = \frac{6\sqrt{2}}{\sqrt{2}} = 6 \text{ cm}.
\]
\[
S = a^2 = 36 \text{ cm}^2, \qquad P = 4 \cdot 6 = 24 \text{ cm}.
\]
\end{example}

\begin{exercise}
Kvadrato plotas $S = 49$ cm$^2$. Raskite jo perimetrą ir įstrižainės ilgį.
\end{exercise}

% -------------------------------------------------------
\subsubsection*{Specialiųjų lygiagretainių palyginimas}

Žemiau pateikiama keturkampių hierarchija ir pagrindinių formulių suvestinė:

\begin{center}
\begin{tabular}{|l|c|c|c|c|}
\hline
\textbf{Savybė} & \textbf{Lygiagretainis} & \textbf{Stačiakampis} & \textbf{Rombas} & \textbf{Kvadratas} \\
\hline
Priešingos kraštinės lygios    & \checkmark & \checkmark & \checkmark & \checkmark \\
Visos kraštinės lygios         &            &            & \checkmark & \checkmark \\
Visi kampai $90°$              &            & \checkmark &            & \checkmark \\
Įstrižainės lygios             &            & \checkmark &            & \checkmark \\
Įstrižainės statmenos          &            &            & \checkmark & \checkmark \\
Įstrižainės dalija pusiau      & \checkmark & \checkmark & \checkmark & \checkmark \\
\hline
\textbf{Plotas}                & $ab\sin\alpha$ & $ab$ & $\dfrac{d_1 d_2}{2}$ & $a^2$ \\
\hline
\textbf{Perimetras}            & $2(a+b)$   & $2(a+b)$   & $4a$       & $4a$  \\
\hline
\end{tabular}
\end{center}

% -------------------------------------------------------
\subsection{Trapecija}
% -------------------------------------------------------

\begin{definition}[Trapecija]
\textbf{Trapecija} – tai keturkampis, kurio \emph{tik} dvi priešingosios kraštinės yra lygiagrečios. Lygiagrečiosios kraštinės vadinamos \textbf{pagrindais} ($a$ ir $b$, kur $a > b$), o kitos dvi – \textbf{šoninėmis kraštinėmis} ($c$ ir $d$).
\end{definition}

\subsubsection*{Trapecijos rūšys}

\begin{enumerate}
    \item \textbf{Lygiašonė trapecija} – šoninės kraštinės lygios: $c = d$. Jos kampai prie pagrindo lygūs, o įstrižainės lygios: $d_1 = d_2$.
    \item \textbf{Stačioji trapecija} – viena šoninė kraštinė statmena pagrindams (ji yra kartu ir aukštinė).
    \item \textbf{Bendrosios rūšies trapecija} – nei lygiašonė, nei stačioji.
\end{enumerate}

\subsubsection*{Trapecijos savybės}

\begin{theorem}[Trapecijos savybės]
\begin{enumerate}
    \item Kampų, esančių prie tos pačios šoninės kraštinės, suma lygi $180°$:
    \[
    \angle A + \angle B = 180°, \quad \angle C + \angle D = 180°.
    \]
    \item \textbf{Vidurio linija} – atkarpa, jungianti šoninių kraštinių vidurius. Ji yra lygiagreti pagrindams ir lygi jų sumos pusei:
    \[
    m = \frac{a + b}{2}.
    \]
    \item Lygiašonės trapecijos įstrižainės lygios: $d_1 = d_2$.
    \item Lygiašonės trapecijos kampai prie to paties pagrindo yra lygūs.
\end{enumerate}
\end{theorem}

\subsubsection*{Trapecijos aukštinė}

Trapecijos \textbf{aukštinė} $h$ – tai statmuo, nuleistas iš vieno pagrindo į kitą.

\begin{remark}
Stačiojoje trapecijoje aukštinė sutampa su viena iš šoninių kraštinių. Lygiašonėje trapecijoje aukštinės pėdos nuo viršūnių yra vienodai nutolusios nuo kraštinių.
\end{remark}

\subsubsection*{Trapecijos plotas ir perimetras}

\begin{tcolorbox}[colback=blue!5!white, colframe=blue!60!black, title=Trapecijos ploto formulės]
\[
S = \frac{a+b}{2} \cdot h = m \cdot h,
\]
kur $a$ ir $b$ – pagrindų ilgiai, $h$ – aukštinė, $m = \dfrac{a+b}{2}$ – vidurio linija.
\end{tcolorbox}

\begin{tcolorbox}[colback=green!5!white, colframe=green!60!black, title=Trapecijos perimetras]
\[
P = a + b + c + d,
\]
kur $a$, $b$ – pagrindai, $c$, $d$ – šoninės kraštinės.
\end{tcolorbox}

\subsubsection*{Lygiašonės trapecijos šoninė kraštinė}

Jeigu lygiašonės trapecijos pagrindai yra $a$ ir $b$ ($a > b$) ir aukštinė $h$, tai šoninę kraštinę $c$ galima rasti iš Pitagoro teoremos:
\[
c = \sqrt{h^2 + \left(\frac{a-b}{2}\right)^2}.
\]

\begin{example}
Trapecijos pagrindai $a = 14$ cm ir $b = 6$ cm, aukštinė $h = 5$ cm. Raskite trapecijos plotą ir vidurio liniją.

\textbf{Sprendimas.}
\[
m = \frac{a+b}{2} = \frac{14+6}{2} = 10 \text{ cm}.
\]
\[
S = m \cdot h = 10 \cdot 5 = 50 \text{ cm}^2.
\]
Arba tiesiogiai:
\[
S = \frac{a+b}{2} \cdot h = \frac{14+6}{2} \cdot 5 = 50 \text{ cm}^2.
\]
\end{example}

\begin{example}
Lygiašonės trapecijos pagrindai $a = 10$ cm ir $b = 4$ cm, šoninė kraštinė $c = 5$ cm. Raskite trapecijos plotą.

\textbf{Sprendimas.} Pirmiausia randame aukštinę. Nuleidę aukštines iš trumpesnio pagrindo galų, gausime du stačiuosius trikampius, kurių vienas statinis yra $\dfrac{a-b}{2} = \dfrac{10-4}{2} = 3$ cm, o įžambinė $c = 5$ cm. Tada:
\[
h = \sqrt{c^2 - \left(\frac{a-b}{2}\right)^2} = \sqrt{25 - 9} = \sqrt{16} = 4 \text{ cm}.
\]
\[
S = \frac{a+b}{2} \cdot h = \frac{10+4}{2} \cdot 4 = 7 \cdot 4 = 28 \text{ cm}^2.
\]
\end{example}

\begin{exercise}
Trapecijos pagrindai $a = 20$ cm ir $b = 12$ cm, o šoninės kraštinės $c = 10$ cm ir $d = 6$ cm. Raskite trapecijos perimetrą. Ar ši trapecija yra lygiašonė?
\end{exercise}

\begin{exercise}
Lygiašonės trapecijos plotas $S = 80$ cm$^2$, vidurio linija $m = 10$ cm. Raskite trapecijos aukštinę.
\end{exercise}

\section{Apskritimas ir skritulys}

\subsection{Apskritimo ir skritulio savybės}

\begin{definition}
\textbf{Apskritimas} – tai plokštumos taškų aibė, nutolusių nuo duotojo taško (centro) vienodu atstumu $r$, vadinamu \textbf{spinduliu}.
\end{definition}

\begin{definition}
\textbf{Skritulys} – tai plokštumos dalis (taškų aibė), apribota apskritimu. Skritulį sudaro visi taškai, nutolę nuo centro $O$ atstumu, ne didesniu už spindulį $r$:
\[
D(O,\,r) = \{\, P \in \text{plokštuma} \mid |OP| \leq r \,\}.
\]
\end{definition}

\begin{remark}
Svarbu skirti sąvokas: \emph{apskritimas} yra tik linija (kreivė), o \emph{skritulys} – visa plokštumos dalis, apribota tuo apskritimu.
\end{remark}

\subsubsection*{Pagrindinės sąvokos}

\begin{itemize}
  \item \textbf{Centras} $O$ – taškas, nuo kurio visi apskritimo taškai yra vienodai nutolę.
  \item \textbf{Spindulys} $r$ – atkarpa, jungianti centrą su bet kuriuo apskritimo tašku; taip pat šios atkarpos ilgis.
  \item \textbf{Skersmuo} $d = 2r$ – apskritimo styga, einanti per centrą.
  \item \textbf{Styga} – atkarpa, jungianti du apskritimo taškus. Ilgiausia styga yra skersmuo.
  \item \textbf{Lankas} – apskritimo dalis tarp dviejų taškų.
  \item \textbf{Centrinis kampas} – kampas, kurio viršūnė yra apskritimo centre.
  \item \textbf{Įbrėžtinis kampas} – kampas, kurio viršūnė yra ant apskritimo, o kraštinės eina per du kitus apskritimo taškus.
\end{itemize}

\subsubsection*{Apskritimo lygtis}

Apskritimo, kurio centras yra taške $O(a,\, b)$ ir spindulys $r$, \textbf{lygtis koordinačių plokštumoje}:
\[
(x - a)^2 + (y - b)^2 = r^2.
\]

Jei centras yra koordinačių pradžioje $O(0,\, 0)$:
\[
x^2 + y^2 = r^2.
\]

\begin{example}
Parašykite apskritimo, kurio centras $O(3,\,-2)$ ir spindulys $r = 5$, lygtį.

\textbf{Sprendimas:}
\[
(x-3)^2 + (y+2)^2 = 25.
\]
\end{example}

\subsubsection*{Apskritimo ilgis ir skritulio plotas}

\begin{tcolorbox}[colback=blue!5!white, colframe=blue!60!black, title=Formulės]
\begin{itemize}
  \item Apskritimo ilgis (perimetras):
  \[
  C = 2\pi r = \pi d.
  \]
  \item Skritulio plotas:
  \[
  S = \pi r^2.
  \]
\end{itemize}
\end{tcolorbox}

\subsubsection*{Lanko ilgis ir išpjovos plotas}

Jei centrinis kampas yra $\alpha$ laipsnių, tai:

\begin{tcolorbox}[colback=green!5!white, colframe=green!60!black, title=Formulės]
\begin{itemize}
  \item Lanko ilgis:
  \[
  l = \frac{\pi r \alpha}{180^{\circ}}.
  \]
  \item Išpjovos plotas:
  \[
  S_{\text{išp}} = \frac{\pi r^2 \alpha}{360^{\circ}}.
  \]
\end{itemize}
\end{tcolorbox}

\begin{example}
Apskaičiuokite lanko, atitinkančio centrinį kampą $\alpha = 60^\circ$, ilgį, jei apskritimo spindulys $r = 6$ cm.

\textbf{Sprendimas:}
\[
l = \frac{\pi \cdot 6 \cdot 60^{\circ}}{180^{\circ}} = \frac{360\pi}{180} = 2\pi \approx 6{,}28 \text{ cm}.
\]
\end{example}

\subsubsection*{Įbrėžtinio kampo teorema}

\begin{theorem}[Įbrėžtinio kampo teorema]
Įbrėžtinis kampas lygus pusei centrinio kampo, paremto tuo pačiu lanku:
\[
\angle \text{įbrėžt.} = \frac{1}{2}\,\angle \text{centrinis}.
\]
\end{theorem}

\begin{corollary}
Visi įbrėžtiniai kampai, paremti tuo pačiu lanku, yra lygūs.
\end{corollary}

\begin{corollary}
Įbrėžtinis kampas, paremtas skersmeniu, yra lygus $90^\circ$ (\textbf{Taleso teorema}).
\end{corollary}

\begin{example}
Apskritimo centre $O$ žinoma, kad centrinis kampas $\angle AOB = 110^\circ$. Raskite įbrėžtinį kampą $\angle ACB$, kur taškas $C$ yra ant mažesnio lanko $AB$ priešingos pusės.

\textbf{Sprendimas:} Taškas $C$ yra ant didesnio lanko, todėl jis remia mažesnį lanką $AB$.
\[
\angle ACB = \frac{\angle AOB}{2} = \frac{110^\circ}{2} = 55^\circ.
\]
\end{example}

\subsubsection*{Liestinė ir spindulys}

\begin{theorem}
Apskritimo spindulys, nubrėžtas į liestinės lietimosi tašką, yra statmenas liestinei:
\[
OT \perp \ell,
\]
kur $T$ – lietimosi taškas, $\ell$ – liestinė.
\end{theorem}

\begin{theorem}
Iš taško, esančio už apskritimo, galima nubrėžti dvi apskritimo liestines, ir jų atkarpų nuo to taško iki lietimosi taškų ilgiai yra lygūs:
\[
|PA| = |PB|,
\]
kur $A$ ir $B$ – lietimosi taškai, $P$ – išorinis taškas.
\end{theorem}

\begin{example}
Iš taško $P$ nubrėžtos dvi apskritimo liestinės. Atstumas nuo $P$ iki centro $O$ yra $10$ cm, o spindulys $r = 6$ cm. Raskite liestinių atkarpų ilgį.

\textbf{Sprendimas:} Kadangi $OA \perp PA$, trikampis $OAP$ yra statusis:
\[
|PA|^2 = |PO|^2 - r^2 = 10^2 - 6^2 = 100 - 36 = 64,
\]
\[
|PA| = 8 \text{ cm}.
\]
\end{example}

% -------------------------------------------------------
\subsection{Apskritimo ir tiesės tarpusavio padėtis}

Tiesė plokštumoje gali turėti tris skirtingas tarpusavio padėtis su apskritimu, priklausomai nuo atstumo $d$ nuo apskritimo centro $O$ iki tiesės.

\begin{tcolorbox}[colback=yellow!10!white, colframe=orange!80!black, title=Tiesės ir apskritimo padėtys]
Tegul $r$ – apskritimo spindulys, $d$ – atstumas nuo centro iki tiesės.
\begin{itemize}
  \item \textbf{Kirstinė} ($d < r$): tiesiė kerta apskritimą dviejuose taškuose.
  \item \textbf{Liestinė} ($d = r$): tiesė liečia apskritimą viename taške.
  \item \textbf{Nesikertanti} ($d > r$): tiesė neturi bendrų taškų su apskritimu.
\end{itemize}
\end{tcolorbox}

\subsubsection*{Atstumas nuo taško iki tiesės}

Atstumas nuo taško $O(x_0,\, y_0)$ iki tiesės $ax + by + c = 0$ apskaičiuojamas formule:
\[
d = \frac{|ax_0 + by_0 + c|}{\sqrt{a^2 + b^2}}.
\]

\begin{example}
Nustatykite tiesės $3x + 4y - 25 = 0$ ir apskritimo $x^2 + y^2 = 25$ tarpusavio padėtį.

\textbf{Sprendimas:} Apskritimo centras $O(0, 0)$, spindulys $r = 5$.

Atstumas nuo $O$ iki tiesės:
\[
d = \frac{|3\cdot 0 + 4\cdot 0 - 25|}{\sqrt{3^2 + 4^2}} = \frac{25}{\sqrt{25}} = \frac{25}{5} = 5.
\]

Kadangi $d = r = 5$, tiesė yra \textbf{liestinė}.
\end{example}

\begin{example}
Nustatykite tiesės $y = x + 1$ ir apskritimo $(x-2)^2 + (y-3)^2 = 9$ tarpusavio padėtį.

\textbf{Sprendimas:} Apskritimo centras $O(2, 3)$, spindulys $r = 3$.

Perrašome tiesę standartiniu pavidalu: $x - y + 1 = 0$.

\[
d = \frac{|1\cdot 2 + (-1)\cdot 3 + 1|}{\sqrt{1^2 + (-1)^2}} = \frac{|2 - 3 + 1|}{\sqrt{2}} = \frac{0}{\sqrt{2}} = 0.
\]

Kadangi $d = 0 < r = 3$, tiesė yra \textbf{kirstinė} (ir eina per centrą – tai skersmuo).
\end{example}

\subsubsection*{Dviejų apskritimų tarpusavio padėtis}

Tegul duoti du apskritimai su centrais $O_1$, $O_2$ ir spinduliais $r_1$, $r_2$ ($r_1 \geq r_2$). Žymime $d = |O_1 O_2|$.

\begin{center}
\begin{tabular}{|l|c|l|}
\hline
\textbf{Padėtis} & \textbf{Sąlyga} & \textbf{Bendrų taškų sk.} \\
\hline
Išoriniai & $d > r_1 + r_2$ & 0 \\
Išoriškai liečiasi & $d = r_1 + r_2$ & 1 \\
Kertasi & $|r_1 - r_2| < d < r_1 + r_2$ & 2 \\
Vidiškai liečiasi & $d = r_1 - r_2$ & 1 \\
Vidinis & $d < r_1 - r_2$ & 0 \\
Sutampa & $d = 0,\ r_1 = r_2$ & $\infty$ \\
\hline
\end{tabular}
\end{center}

\begin{example}
Ar apskritimai $(x-1)^2 + (y-1)^2 = 9$ ir $(x-5)^2 + (y-4)^2 = 4$ kertasi?

\textbf{Sprendimas:} $O_1(1,1)$, $r_1 = 3$; $O_2(5,4)$, $r_2 = 2$.
\[
d = \sqrt{(5-1)^2 + (4-1)^2} = \sqrt{16 + 9} = \sqrt{25} = 5.
\]
Tikriname: $r_1 + r_2 = 5$ ir $d = 5$, todėl apskritimai \textbf{išoriškai liečiasi} viename taške.
\end{example}

% -------------------------------------------------------
\subsection{Įbrėžtinis ir apibrėžtinis apskritimas}

\subsubsection*{Apibrėžtinis apskritimas (apie trikampį)}

\begin{definition}
\textbf{Apibrėžtinis apskritimas} – apskritimas, einantis per visas trikampio (ar kito daugiakampio) viršūnes.
\end{definition}

\begin{theorem}
Apie kiekvieną trikampį galima apibrėžti apskritimą, ir jo centras yra trikampio \textbf{kraštinių vidurio statmenų susikirtimo taškas} (vadinamas \emph{apcentru}).
\end{theorem}

\begin{tcolorbox}[colback=blue!5!white, colframe=blue!60!black, title=Apibrėžtinio apskritimo spindulys]
Apie trikampį $ABC$ apibrėžto apskritimo spindulys apskaičiuojamas pagal \textbf{sinusų teoremą}:
\[
R = \frac{a}{2\sin\alpha} = \frac{b}{2\sin\beta} = \frac{c}{2\sin\gamma},
\]
kur $a, b, c$ – trikampio kraštinės, $\alpha, \beta, \gamma$ – priešingieji kampai.

Taip pat naudojama formulė per trikampio plotą $S$:
\[
R = \frac{abc}{4S}.
\]
\end{tcolorbox}

\subsubsection*{Apibrėžtinis apskritimas ypatingiems trikampiams}

\begin{itemize}
  \item \textbf{Lygiakraštis trikampis} (kraštinė $a$):
  \[
  R = \frac{a}{\sqrt{3}} = \frac{a\sqrt{3}}{3}.
  \]
  \item \textbf{Stačiasis trikampis} (įžambinė $c$):
  \[
  R = \frac{c}{2}.
  \]
  (Apibrėžtinio apskritimo centras yra įžambinės vidurys.)
  \item \textbf{Lygiašonis trikampis} su kraštinėmis $a, a, b$:
  \[
  R = \frac{a^2 b}{\sqrt{(2a+b)(2a-b)}\cdot b} = \frac{a^2}{\sqrt{4a^2 - b^2}}.
  \]
\end{itemize}

\subsubsection*{Įbrėžtinis apskritimas (į trikampį)}

\begin{definition}
\textbf{Įbrėžtinis apskritimas} – apskritimas, liečiantis visas trikampio (ar kito daugiakampio) kraštines iš vidaus.
\end{definition}

\begin{theorem}
Į kiekvieną trikampį galima įbrėžti apskritimą, ir jo centras yra trikampio \textbf{pusiaukampinių susikirtimo taškas} (vadinamas \emph{incentru}).
\end{theorem}

\begin{tcolorbox}[colback=green!5!white, colframe=green!60!black, title=Įbrėžtinio apskritimo spindulys]
Į trikampį $ABC$ įbrėžto apskritimo spindulys:
\[
r = \frac{S}{p},
\]
kur $S$ – trikampio plotas, $p = \dfrac{a+b+c}{2}$ – trikampio pusperimetris.
\end{tcolorbox}

\subsubsection*{Įbrėžtinis apskritimas ypatingiems trikampiams}

\begin{itemize}
  \item \textbf{Lygiakraštis trikampis} (kraštinė $a$):
  \[
  r = \frac{a\sqrt{3}}{6} = \frac{a}{2\sqrt{3}}.
  \]
  \item \textbf{Stačiasis trikampis} (statiniai $a, b$, įžambinė $c$):
  \[
  r = \frac{a + b - c}{2}.
  \]
  \item \textbf{Pastaba:} Lygiakraščio trikampio atveju $R = 2r$.
\end{itemize}

\begin{example}
Trikampio kraštinės yra $a = 5$, $b = 12$, $c = 13$ cm. Ar trikampis stačiasis? Apskaičiuokite įbrėžtinio ir apibrėžtinio apskritimų spindulius.

\textbf{Sprendimas:}

\textbf{1.} Tikriname: $5^2 + 12^2 = 25 + 144 = 169 = 13^2$. Taip, trikampis \textbf{stačiasis}.

\textbf{2.} Trikampio plotas: $S = \dfrac{1}{2} \cdot 5 \cdot 12 = 30$ cm\textsuperscript{2}.

\textbf{3.} Pusperimetris: $p = \dfrac{5 + 12 + 13}{2} = 15$ cm.

\textbf{4.} Įbrėžtinio apskritimo spindulys:
\[
r = \frac{S}{p} = \frac{30}{15} = 2 \text{ cm}.
\]

\textbf{5.} Apibrėžtinio apskritimo spindulys (stačiajam trikampiui):
\[
R = \frac{c}{2} = \frac{13}{2} = 6{,}5 \text{ cm}.
\]
\end{example}

\subsubsection*{Įbrėžtinis ir apibrėžtinis apskritimas keturkampiuose}

\begin{theorem}
Į iškiląjį keturkampį galima įbrėžti apskritimą \textbf{tada ir tik tada}, kai jo priešingų kraštinių sumos yra lygios:
\[
a + c = b + d,
\]
kur $a, b, c, d$ – keturkampio kraštinės iš eilės.
\end{theorem}

\begin{theorem}
Apie iškiląjį keturkampį galima apibrėžti apskritimą \textbf{tada ir tik tada}, kai jo priešingų kampų suma lygi $180^\circ$:
\[
\alpha + \gamma = \beta + \delta = 180^\circ.
\]
Toks keturkampis vadinamas \textbf{cikliniu}.
\end{theorem}

\begin{example}
Stačiakampis keturkampis $ABCD$ turi kraštines $AB = 3$, $BC = 5$, $CD = 7$, $DA = x$. Raskite $x$, jei apskritimas galėtų būti įbrėžtas.

\textbf{Sprendimas:}
\[
AB + CD = BC + DA \implies 3 + 7 = 5 + x \implies x = 5.
\]
\end{example}

\begin{exercise}
Trikampio plotas lygus $84$ cm\textsuperscript{2}, o kraštinės yra $13$, $14$ ir $15$ cm. Apskaičiuokite į trikampį įbrėžto apskritimo spindulį.
\end{exercise}

\begin{exercise}
Apskritimo, kurio centras $O(-1, 2)$ ir spindulys $r = 4$, atžvilgiu nustatykite tiesės $3x - 4y + 15 = 0$ padėtį.
\end{exercise}

\begin{exercise}
Įrodykite, kad taškas $P(3, 4)$ yra ant apskritimo $x^2 + y^2 = 25$, ir nubrėžkite liestinę šiame taške. Raskite tos liestinės lygtį.
\end{exercise}

\section{Koordinačių sistema plokštumoje}

Koordinačių sistema plokštumoje – tai sistema, leidžianti tiksliai apibūdinti kiekvieno taško padėtį plokštumoje naudojant du skaičius (koordinates). Ji sudaroma iš dviejų statmenų koordinačių ašių: horizontaliosios \textbf{abscisių ašies} $Ox$ ir vertikaliosios \textbf{ordinačių ašies} $Oy$. Šių ašių susikirtimo taškas vadinamas \textbf{koordinačių pradžia} $O(0;0)$.

\begin{definition}[Taško koordinatės]
Taškas $A$ plokštumoje aprašomas pora $(x; y)$, kur $x$ – taško \textbf{abscisė} (atstumas iki $Oy$ ašies), o $y$ – taško \textbf{ordinatė} (atstumas iki $Ox$ ašies). Rašome $A(x; y)$.
\end{definition}

Koordinačių ašys plokštumą dalija į keturis \textbf{ketvirčius}:
\begin{itemize}
    \item \textbf{I ketvirtis}: $x > 0$, $y > 0$
    \item \textbf{II ketvirtis}: $x < 0$, $y > 0$
    \item \textbf{III ketvirtis}: $x < 0$, $y < 0$
    \item \textbf{IV ketvirtis}: $x > 0$, $y < 0$
\end{itemize}

% ─────────────────────────────────────────────
\subsection{Atstumo formulė}
% ─────────────────────────────────────────────

\begin{definition}[Atstumas tarp dviejų taškų]
Atstumas tarp dviejų koordinačių plokštumos taškų $A(x_1;\, y_1)$ ir $B(x_2;\, y_2)$ lygus:
\[
    |AB| = \sqrt{(x_2 - x_1)^2 + (y_2 - y_1)^2}.
\]
\end{definition}

\textbf{Įrodymo idėja.} Taškus $A$ ir $B$ sujunkime atkarpa. Statykime iš $A$ statmeną liniją $Ox$ krypčiai ir iš $B$ – liniją $Oy$ krypčiai. Gauname stačiojo trikampio $AKB$ (kur $K(x_2; y_1)$) katetus:
\[
    |AK| = |x_2 - x_1|, \qquad |KB| = |y_2 - y_1|.
\]
Pagal Pitagoro teoremą:
\[
    |AB|^2 = |AK|^2 + |KB|^2 = (x_2 - x_1)^2 + (y_2 - y_1)^2,
\]
iš kur gauname atstumo formulę. $\square$

\begin{remark}
Specialieji atvejai:
\begin{itemize}
    \item Jei $A$ ir $B$ ant $Ox$ ašies ($y_1 = y_2 = 0$): $\;|AB| = |x_2 - x_1|$.
    \item Jei $A$ ir $B$ ant $Oy$ ašies ($x_1 = x_2 = 0$): $\;|AB| = |y_2 - y_1|$.
    \item Atstumas nuo taško $A(x_1; y_1)$ iki koordinačių pradžios $O(0;0)$:
    \[
        |OA| = \sqrt{x_1^2 + y_1^2}.
    \]
\end{itemize}
\end{remark}

\begin{example}
Raskite atstumą tarp taškų $A(2;\, 3)$ ir $B(-2;\, 1)$.

\textbf{Sprendimas:}
\[
    |AB| = \sqrt{(-2-2)^2 + (1-3)^2} = \sqrt{(-4)^2 + (-2)^2} = \sqrt{16 + 4} = \sqrt{20} = 2\sqrt{5}.
\]
\end{example}

\begin{example}
Įrodykite, kad taškai $A(0;\,0)$, $B(4;\,0)$ ir $C(2;\,2\sqrt{3})$ sudaro lygiakraštį trikampį.

\textbf{Sprendimas:}
\begin{align*}
    |AB| &= \sqrt{(4-0)^2 + (0-0)^2} = 4, \\
    |BC| &= \sqrt{(2-4)^2 + (2\sqrt{3}-0)^2} = \sqrt{4 + 12} = 4, \\
    |CA| &= \sqrt{(0-2)^2 + (0-2\sqrt{3})^2} = \sqrt{4 + 12} = 4.
\end{align*}
Kadangi $|AB| = |BC| = |CA| = 4$, trikampis yra lygiakraštis. $\square$
\end{example}

\begin{exercise}
Raskite atstumus tarp šių taškų porų:
\begin{enumerate}[label=\alph*)]
    \item $P(1;\,2)$ ir $Q(4;\,6)$
    \item $M(-3;\,5)$ ir $N(2;\,-7)$
    \item $A(-1;\,-1)$ ir $B(3;\,2)$
\end{enumerate}
\end{exercise}

% ─────────────────────────────────────────────
\subsection{Atkarpos vidurio formulė}
% ─────────────────────────────────────────────

\begin{definition}[Atkarpos vidurio taškas]
Atkarpos $AB$, kur $A(x_1;\, y_1)$ ir $B(x_2;\, y_2)$, vidurio taško $V$ koordinatės:
\[
    V = \left(\frac{x_1 + x_2}{2};\; \frac{y_1 + y_2}{2}\right).
\]
\end{definition}

\textbf{Įrodymo idėja.} Taškas $V(x_v;\, y_v)$ yra atkarpos $AB$ vidurys, kai $|AV| = |VB|$. Kadangi koordinatės kinta tolygiai palei atkarpą, vidurio koordinatė yra kraštinių koordinačių aritmetinis vidurkis. $\square$

\begin{example}
Raskite atkarpos $AB$ vidurio taško koordinates, jei $A(-4;\, 6)$ ir $B(2;\, -2)$.

\textbf{Sprendimas:}
\[
    V = \left(\frac{-4+2}{2};\; \frac{6+(-2)}{2}\right) = \left(\frac{-2}{2};\; \frac{4}{2}\right) = (-1;\; 2).
\]
\end{example}

\begin{example}
Žinoma, kad taškas $V(3;\, 1)$ yra atkarpos $AB$ vidurys. Raskite taško $B$ koordinates, jei $A(1;\, -3)$.

\textbf{Sprendimas.} Iš vidurio formulės:
\[
    \frac{x_1 + x_2}{2} = 3 \implies x_2 = 2 \cdot 3 - 1 = 5,
\]
\[
    \frac{y_1 + y_2}{2} = 1 \implies y_2 = 2 \cdot 1 - (-3) = 5.
\]
Taigi $B(5;\, 5)$.
\end{example}

\begin{remark}
Vidurio formulė naudojama ir koordinačių geometrijoje, sprendžiant uždavinius apie:
\begin{itemize}
    \item trikampio medianas (jos eina per viršūnę ir priešingosios kraštinės vidurį),
    \item lygiagretainio savybes (įstrižainės pusiaudala viena kitą),
    \item apskritimo centro radimą pagal skersmens galus.
\end{itemize}
\end{remark}

\begin{exercise}
\begin{enumerate}[label=\alph*)]
    \item Raskite atkarpos $MN$ vidurio taško koordinates, jei $M(5;\,-3)$ ir $N(-1;\,7)$.
    \item Žinoma, kad atkarpos $PQ$ vidurio taškas yra $S(0;\,4)$ ir $P(-2;\,1)$. Raskite $Q$ koordinates.
    \item Lygiagretainio $ABCD$ trys viršūnės yra $A(1;\,2)$, $B(5;\,0)$, $C(6;\,4)$. Raskite $D$ koordinates.
\end{enumerate}
\end{exercise}

% ─────────────────────────────────────────────
\subsection{Tiesės lygtis plokštumoje}
% ─────────────────────────────────────────────

Tiesė plokštumoje gali būti aprašoma keliais lygiaverčiais būdais. Svarbu mokėti atpažinti ir taikyti kiekvieną.

% ── 1. Bendroji lygtis ──────────────────────
\subsubsection*{Bendroji tiesės lygtis}

\begin{definition}[Bendroji forma]
Kiekviena tiesė plokštumoje gali būti užrašyta \textbf{bendrosios lygties} pavidalu:
\[
    Ax + By + C = 0,
\]
kur $A$, $B$, $C \in \mathbb{R}$ ir $(A, B) \neq (0, 0)$.
\end{definition}

Koeficientai turi aiškią geometrinę prasmę: vektorius $\vec{n} = (A;\, B)$ yra tiesės \textbf{normalės vektorius} (statmenas tiesės krypčiai).

% ── 2. Lygtis su nuolydžio koeficientu ────
\subsubsection*{Lygtis su nuolydžio koeficientu}

\begin{definition}[Nuolydžio koeficiento forma]
Jei tiesė nėra lygiagreti $Oy$ ašiai ($B \neq 0$), ją galima užrašyti:
\[
    y = kx + b,
\]
kur:
\begin{itemize}
    \item $k$ – \textbf{nuolydžio koeficientas} (charakterizuoja tiesės pasvirumą),
    \item $b$ – \textbf{ordinačių ašies atkarpėlė} (tiesė kerta $Oy$ taške $(0;\, b)$).
\end{itemize}
\end{definition}

\begin{definition}[Nuolydžio koeficiento geometrinė prasmė]
Jei $\alpha$ – kampas tarp tiesės ir teigiamosios $Ox$ ašies pusės, tai:
\[
    k = \tg \alpha.
\]
\begin{itemize}
    \item $k > 0$: tiesė kyla iš kairės į dešinę.
    \item $k < 0$: tiesė leidžiasi iš kairės į dešinę.
    \item $k = 0$: tiesė lygiagreti $Ox$ ašiai.
\end{itemize}
\end{definition}

Ryšys tarp bendrosios lygties ir nuolydžio koeficiento:
\[
    k = -\frac{A}{B}, \qquad b = -\frac{C}{B} \quad (B \neq 0).
\]

% ── 3. Tiesė per tašką su nuolydžiu ────────
\subsubsection*{Tiesės lygtis, einanti per tašką su duotu nuolydžiu}

\begin{theorem}
Tiesės, einančios per tašką $M_0(x_0;\, y_0)$ su nuolydžio koeficientu $k$, lygtis:
\[
    y - y_0 = k(x - x_0).
\]
\end{theorem}

\begin{example}
Užrašykite tiesės lygtį, einančią per tašką $A(2;\, -1)$ su nuolydžio koeficientu $k = 3$.

\textbf{Sprendimas:}
\[
    y - (-1) = 3(x - 2) \implies y + 1 = 3x - 6 \implies y = 3x - 7.
\]
\end{example}

% ── 4. Tiesė per du taškus ──────────────────
\subsubsection*{Tiesės lygtis, einanti per du taškus}

\begin{theorem}
Tiesės, einančios per du skirtingus taškus $A(x_1;\, y_1)$ ir $B(x_2;\, y_2)$ ($x_1 \neq x_2$), lygtis:
\[
    \frac{y - y_1}{y_2 - y_1} = \frac{x - x_1}{x_2 - x_1},
\]
arba ekvivalenčiai nuolydžio koeficientas:
\[
    k = \frac{y_2 - y_1}{x_2 - x_1}.
\]
Kai $x_1 = x_2$ (tiesė lygiagreti $Oy$), tiesės lygtis yra $x = x_1$.
\end{theorem}

\begin{example}
Raskite tiesės lygtį, einančią per taškus $A(1;\, 2)$ ir $B(3;\, 8)$.

\textbf{Sprendimas.} Pirmiausia randame nuolydžio koeficientą:
\[
    k = \frac{8 - 2}{3 - 1} = \frac{6}{2} = 3.
\]
Tada taikome formulę per tašką $A(1;\,2)$:
\[
    y - 2 = 3(x - 1) \implies y = 3x - 1.
\]
\end{example}

% ── 5. Lygiagrečios ir statmenos tiesės ────
\subsubsection*{Tiesių tarpusavio padėtis}

\begin{theorem}[Lygiagretumo sąlyga]
Dvi tiesės $y = k_1 x + b_1$ ir $y = k_2 x + b_2$ yra \textbf{lygiagrečios}, jei ir tik jei:
\[
    k_1 = k_2 \quad \text{ir} \quad b_1 \neq b_2.
\]
Jei $k_1 = k_2$ ir $b_1 = b_2$ – tiesės sutampa.
\end{theorem}

\begin{theorem}[Statmenumo sąlyga]
Dvi tiesės $y = k_1 x + b_1$ ir $y = k_2 x + b_2$ yra \textbf{statmenos}, jei ir tik jei:
\[
    k_1 \cdot k_2 = -1.
\]
\end{theorem}

\begin{example}
Parašykite tiesės, einančios per tašką $P(4;\,1)$ ir statmenos tiesei $y = 2x + 3$, lygtį.

\textbf{Sprendimas.} Duotos tiesės $k_1 = 2$. Statmenos tiesės nuolydžio koeficientas:
\[
    k_2 = -\frac{1}{k_1} = -\frac{1}{2}.
\]
Tiesės lygtis per tašką $P(4;\,1)$:
\[
    y - 1 = -\frac{1}{2}(x - 4) \implies y = -\frac{1}{2}x + 3.
\]
\end{example}

% ── 6. Atstumas nuo taško iki tiesės ───────
\subsubsection*{Atstumas nuo taško iki tiesės}

\begin{theorem}[Atstumas nuo taško iki tiesės]
Atstumas $d$ nuo taško $M_0(x_0;\, y_0)$ iki tiesės $Ax + By + C = 0$ lygus:
\[
    d = \frac{|A x_0 + B y_0 + C|}{\sqrt{A^2 + B^2}}.
\]
\end{theorem}

\begin{example}
Raskite atstumą nuo taško $M(3;\, -2)$ iki tiesės $3x - 4y + 5 = 0$.

\textbf{Sprendimas:}
\[
    d = \frac{|3 \cdot 3 + (-4)(-2) + 5|}{\sqrt{3^2 + (-4)^2}} = \frac{|9 + 8 + 5|}{\sqrt{9 + 16}} = \frac{22}{5} = 4{,}4.
\]
\end{example}

\begin{remark}
Atstumo nuo taško iki tiesės formulė dažnai naudojama:
\begin{itemize}
    \item apskaičiuojant trikampio aukštinę pagal viršūnę ir priešingąją kraštinę,
    \item nustatant apskritimo ir tiesės tarpusavio padėtį (lyginant $d$ su spinduliu $r$),
    \item sprendžiant optimizavimo uždavinius.
\end{itemize}
\end{remark}

\begin{exercise}
\begin{enumerate}[label=\alph*)]
    \item Parašykite tiesės, einančios per taškus $A(-1;\,3)$ ir $B(2;\,-3)$, lygtį.
    \item Ar tiesės $y = 3x - 1$ ir $6x - 2y + 4 = 0$ yra lygiagrečios? Pagrįskite.
    \item Raskite tiesės, einančios per tašką $(2;\,5)$ ir lygiagrečios $Ox$ ašiai, lygtį.
    \item Raskite atstumą nuo koordinačių pradžios $O(0;\,0)$ iki tiesės $5x + 12y - 26 = 0$.
    \item Parašykite tiesės, statmenos tiesei $2x - y + 1 = 0$ ir einančios per tašką $(4;\,3)$, lygtį.
\end{enumerate}
\end{exercise}

\section{Plotai ir perimetrai}

\subsection{Pagrindinės sąvokos}

\begin{definition}
\textbf{Perimetras} – uždaros plokščios figūros ribos (kontūro) ilgis. Žymimas raide $P$.
\end{definition}

\begin{definition}
\textbf{Plotas} – dydis, apibūdinantis plokščios figūros dydį plokštumoje. Žymimas raide $S$.
\end{definition}

Plotų matavimo vienetai: $\text{mm}^2$, $\text{cm}^2$, $\text{dm}^2$, $\text{m}^2$, $\text{km}^2$. \\
Vienetų perskaičiavimas: $1\,\text{m}^2 = 100\,\text{dm}^2 = 10\,000\,\text{cm}^2 = 1\,000\,000\,\text{mm}^2$.

% -------------------------------------------------------
\subsection{Trikampis}

\begin{tcolorbox}[title={Trikampio perimetras}]
\[
P = a + b + c
\]
čia $a, b, c$ – trikampio kraštinių ilgiai.
\end{tcolorbox}

\subsubsection*{Trikampio ploto formulės}

Egzistuoja kelios lygiavertės trikampio ploto formulės:

\begin{enumerate}
    \item \textbf{Pagal kraštinę ir aukštinę:}
    \[
    S = \frac{1}{2} \cdot a \cdot h_a
    \]
    čia $a$ – kraštinė, $h_a$ – į ją nubrėžta aukštinė.

    \item \textbf{Pagal dvi kraštines ir tarp jų esantį kampą (trigonometrinė formulė):}
    \[
    S = \frac{1}{2} \cdot a \cdot b \cdot \sin C
    \]
    čia $a$, $b$ – dvi kraštinės, $C$ – kampas tarp jų.

    \item \textbf{Herono formulė} (pagal tris kraštines):
    \[
    S = \sqrt{p(p-a)(p-b)(p-c)}
    \]
    čia $p = \dfrac{a+b+c}{2}$ – pusperimetris.

    \item \textbf{Pagal pusperimetrį ir įbrėžto apskritimo spindulį:}
    \[
    S = p \cdot r
    \]
    čia $r$ – įbrėžto apskritimo spindulys.

    \item \textbf{Pagal kraštines ir apibrėžto apskritimo spindulį:}
    \[
    S = \frac{a \cdot b \cdot c}{4R}
    \]
    čia $R$ – apibrėžto apskritimo spindulys.
\end{enumerate}

\begin{tcolorbox}[title={VBE formulyno trikampio suvestinė}]
\[
S = \frac{1}{2}ab\sin C = \sqrt{p(p-a)(p-b)(p-c)} = rp = \frac{abc}{4R}
\]
\end{tcolorbox}

\subsubsection*{Specialieji trikampiai}

\begin{itemize}
    \item \textbf{Stačiasis trikampis} (statiniai $a$ ir $b$):
    \[
    S = \frac{1}{2} \cdot a \cdot b, \qquad P = a + b + c
    \]

    \item \textbf{Lygiakraštis trikampis} (kraštinė $a$):
    \[
    S = \frac{a^2\sqrt{3}}{4}, \qquad h = \frac{a\sqrt{3}}{2}, \qquad P = 3a
    \]
\end{itemize}

\begin{example}
Trikampio kraštinės yra $a = 5$, $b = 7$, $c = 8$. Raskite jo plotą.

\textbf{Sprendimas:}
\[
p = \frac{5+7+8}{2} = 10
\]
\[
S = \sqrt{10(10-5)(10-7)(10-8)} = \sqrt{10 \cdot 5 \cdot 3 \cdot 2} = \sqrt{300} = 10\sqrt{3}
\]
\end{example}

% -------------------------------------------------------
\subsection{Keturkampiai}

\subsubsection*{Kvadratas}

\begin{tcolorbox}[title={Kvadratas (kraštinė $a$, įstrižainė $d$)}]
\[
P = 4a, \qquad S = a^2 = \frac{d^2}{2}
\]
\end{tcolorbox}

Papildomos formulės pagal įbrėžtą ir apibrėžtą apskritimą:
\[
S = 2R^2 = 4r^2
\]
čia $R$ – apibrėžto apskritimo, $r$ – įbrėžto apskritimo spindulys.

\subsubsection*{Stačiakampis}

\begin{tcolorbox}[title={Stačiakampis (kraštinės $a$, $b$, įstrižainė $d$)}]
\[
P = 2(a+b), \qquad S = a \cdot b
\]
\[
d = \sqrt{a^2 + b^2}
\]
\end{tcolorbox}

\subsubsection*{Lygiagretainis}

\begin{tcolorbox}[title={Lygiagretainis}]
\[
P = 2(a+b), \qquad S = a \cdot h = a \cdot b \cdot \sin\alpha
\]
\[
S = \frac{1}{2} d_1 d_2 \sin\varphi
\]
čia $h$ – aukštinė į kraštinę $a$, $\alpha$ – kampas tarp kraštinių, \\
$d_1, d_2$ – įstrižainės, $\varphi$ – kampas tarp įstrižainių.
\end{tcolorbox}

\subsubsection*{Rombas}

\begin{tcolorbox}[title={Rombas (kraštinė $a$, įstrižainės $d_1$, $d_2$)}]
\[
P = 4a, \qquad S = a^2\sin\alpha = \frac{d_1 \cdot d_2}{2}
\]
čia $\alpha$ – kampas tarp kraštinių.
\end{tcolorbox}

\begin{remark}
Rombo įstrižainės statmenos viena kitai, todėl visada galioja: $a^2 = \left(\dfrac{d_1}{2}\right)^2 + \left(\dfrac{d_2}{2}\right)^2$.
\end{remark}

\subsubsection*{Trapecija}

\begin{tcolorbox}[title={Trapecija (pagrindai $a$, $b$; aukštinė $h$; vidurinė linija $m$)}]
\[
P = a + b + c + d, \qquad S = \frac{(a+b) \cdot h}{2} = m \cdot h
\]
čia $m = \dfrac{a+b}{2}$ – vidurinė linija, $c$, $d$ – šonai.
\end{tcolorbox}

\subsubsection*{Įbrėžtinis keturkampis}

Keturkampio, į kurį galima įbrėžti apskritimą, plotas:
\[
S = \sqrt{(p-a)(p-b)(p-c)(p-d)}
\]
čia $a, b, c, d$ – kraštinių ilgiai, $p$ – pusperimetris. Tai \textbf{Brahmagupto formulė}.

\begin{example}
Trapecijos pagrindai yra $a = 10$ cm ir $b = 4$ cm, šonai $c = d = 5$ cm. Raskite aukštinę ir plotą.

\textbf{Sprendimas:}

Pažymime $h$ – aukštinę. Kiekviena šono projekcija ant pagrindo:
\[
x = \frac{a - b}{2} = \frac{10 - 4}{2} = 3 \text{ cm}
\]
\[
h = \sqrt{c^2 - x^2} = \sqrt{25 - 9} = \sqrt{16} = 4 \text{ cm}
\]
\[
S = \frac{(10+4) \cdot 4}{2} = 28 \text{ cm}^2
\]
\end{example}

% -------------------------------------------------------
\subsection{Taisyklingasis daugiakampis}

\begin{definition}
\textbf{Taisyklingasis daugiakampis} – daugiakampis, kurio visos kraštinės lygios ir visi kampai lygūs.
\end{definition}

\begin{tcolorbox}[title={Taisyklingasis $n$-kampis (kraštinė $a$)}]
\[
P = n \cdot a
\]
\[
S = \frac{n \cdot a \cdot r}{2} = \frac{1}{2} R^2 n \sin\!\left(\frac{360°}{n}\right)
\]
čia $r$ – įbrėžto apskritimo spindulys (apotema), $R$ – apibrėžto apskritimo spindulys.
\end{tcolorbox}

\subsubsection*{Taisyklingasis šešiakampis}

Taisyklingajame šešiakampyje kraštinė lygi apibrėžto apskritimo spinduliui ($a = R$), todėl:
\[
S = \frac{3\sqrt{3}}{2}\,a^2, \qquad r = \frac{a\sqrt{3}}{2}, \qquad P = 6a
\]

% -------------------------------------------------------
\subsection{Apskritimas ir skritulys}

\begin{definition}
\textbf{Apskritimas} – visų plokštumos taškų, nutolusių nuo duoto centro taško atstumu $R$, aibė. \\
\textbf{Skritulys} – visų plokštumos taškų, nutolusių nuo duoto centro taško atstumu, ne didesniu nei $R$, aibė.
\end{definition}

\begin{tcolorbox}[title={Apskritimas ir skritulys (spindulys $R$)}]
\[
C = 2\pi R \quad \text{(apskritimo ilgis)}, \qquad S = \pi R^2 \quad \text{(skritulio plotas)}
\]
\end{tcolorbox}

\subsubsection*{Lanko ilgis ir išpjovos plotas}

\begin{tcolorbox}[title={Lankas ir išpjova (centrinis kampas $\alpha$ laipsniais)}]
\[
l = \frac{\pi R \alpha}{180°} = \frac{2\pi R \alpha}{360°}
\]
\[
S_{\text{išpj.}} = \frac{\pi R^2 \alpha}{360°} = \frac{R \cdot l}{2}
\]
\end{tcolorbox}

\subsubsection*{Nuopjovos plotas}

Nuopjova – figūra tarp styginės $AB$ ir lanko $AB$. Centrinis kampas $\alpha$ laipsniais:

\[
S_{\text{nuopj.}} = \frac{R^2}{2}\left(\frac{\pi\alpha}{180°} - \sin\alpha\right)
\]

Arba naudojant išpjovą:
\[
S_{\text{nuopj. (mažoji)}} = S_{\text{išpjova}} - S_{\triangle AOB}
\]
čia $O$ – apskritimo centras, $A$, $B$ – styginės galai.

\begin{example}
Apskritimo spindulys $R = 6$ cm, centrinis kampas $\alpha = 60°$. Raskite lanko ilgį ir išpjovos plotą.

\textbf{Sprendimas:}
\[
l = \frac{\pi \cdot 6 \cdot 60°}{180°} = 2\pi \approx 6{,}28 \text{ cm}
\]
\[
S = \frac{\pi \cdot 36 \cdot 60°}{360°} = 6\pi \approx 18{,}85 \text{ cm}^2
\]
\end{example}

% -------------------------------------------------------
\subsection{Ploto skaičiavimas su integralu}

Naudojant apibrėžtinį integralą, galima apskaičiuoti bet kokios figūros, ribojamos kreivėmis, plotą.

\begin{theorem}
Jei $f(x) \geq 0$ intervale $[a;\,b]$, tai figūros, ribotos kreive $y = f(x)$, $x$ ašimi ir tiesėmis $x = a$, $x = b$, plotas:
\[
S = \int_a^b f(x)\,dx
\]
\end{theorem}

\begin{theorem}
Jei $f(x) \leq 0$ intervale $[a;\,b]$:
\[
S = -\int_a^b f(x)\,dx = \int_a^b |f(x)|\,dx
\]
\end{theorem}

\begin{theorem}
Figūros tarp dviejų kreivių $y = f(x)$ (viršutinė) ir $y = g(x)$ (apatinė) plotas intervale $[a;\,b]$:
\[
S = \int_a^b \bigl[f(x) - g(x)\bigr]\,dx
\]
\end{theorem}

\begin{example}
Raskite figūros, ribotos parabolės $y = x^2$ ir tiesės $y = x + 2$, plotą.

\textbf{Sprendimas:}

Pirmiausia randame susikirtimo taškus:
\[
x^2 = x + 2 \implies x^2 - x - 2 = 0 \implies (x-2)(x+1) = 0
\]
\[
x_1 = -1, \quad x_2 = 2
\]

Intervale $[-1;\,2]$ tiesė yra viršuje: $x + 2 \geq x^2$, todėl:
\[
S = \int_{-1}^{2} \bigl[(x+2) - x^2\bigr]\,dx
= \left[\frac{x^2}{2} + 2x - \frac{x^3}{3}\right]_{-1}^{2}
\]
\[
= \left(2 + 4 - \frac{8}{3}\right) - \left(\frac{1}{2} - 2 + \frac{1}{3}\right)
= \frac{10}{3} - \left(-\frac{7}{6}\right) = \frac{20}{6} + \frac{7}{6} = \frac{27}{6} = \frac{9}{2}
\]
\end{example}

% -------------------------------------------------------
\subsection{Formulių suvestinė}

\begin{center}
\renewcommand{\arraystretch}{1.6}
\begin{tabular}{|l|c|c|}
\hline
\textbf{Figūra} & \textbf{Perimetras} $P$ & \textbf{Plotas} $S$ \\
\hline
Trikampis & $a+b+c$ & $\dfrac{1}{2}ah_a$ \\
\hline
Lygiakraštis trikampis & $3a$ & $\dfrac{a^2\sqrt{3}}{4}$ \\
\hline
Kvadratas & $4a$ & $a^2$ \\
\hline
Stačiakampis & $2(a+b)$ & $ab$ \\
\hline
Lygiagretainis & $2(a+b)$ & $ah$ \\
\hline
Rombas & $4a$ & $\dfrac{d_1 d_2}{2}$ \\
\hline
Trapecija & $a+b+c+d$ & $\dfrac{(a+b)h}{2}$ \\
\hline
Apskritimas/Skritulys & $2\pi R$ & $\pi R^2$ \\
\hline
Išpjova & $2R + l$ & $\dfrac{\pi R^2 \alpha}{360°}$ \\
\hline
\end{tabular}
\end{center}

% -------------------------------------------------------
\subsection{Uždaviniai}

\begin{exercise}
Trikampio dvi kraštinės yra 6 cm ir 8 cm, o kampas tarp jų lygus $30°$. Raskite trikampio plotą ir perimetrą.
\end{exercise}

\begin{exercise}
Rombo perimetras yra 40 cm, o viena įstrižainė lygi 12 cm. Raskite rombo plotą.
\end{exercise}

\begin{exercise}
Apskritimo spindulys $R = 9$ cm. Apskaičiuokite išpjovos, atitinkančios $\alpha = 120°$ centrinį kampą, plotą ir lanko ilgį.
\end{exercise}

\begin{exercise}
Raskite figūros, ribotos kreivės $y = \sqrt{x}$, $x$ ašies ir tiesės $x = 4$, plotą.
\end{exercise}

\begin{exercise}
(VBE tipo uždavinys) Lygiašonio trikampio pagrindas yra 10 cm, o šonų ilgiai yra 13 cm. Raskite trikampio plotą, aukštinę į pagrindą, įbrėžto apskritimo spindulį ir apibrėžto apskritimo spindulį.
\end{exercise}

\begin{exercise}
Raskite figūros, ribotos kreivių $y = x^2 - 4$ ir $y = -x^2 + 4$, plotą.
\end{exercise}